\documentclass[11pt]{article}
\usepackage[T1]{fontenc}
\usepackage[margin=1in]{geometry}
\usepackage{amsmath}
\usepackage{amssymb}
\usepackage{array}
\usepackage{longtable}
\usepackage{textcomp}
\usepackage{xurl}
\usepackage[hidelinks]{hyperref}

\newcommand{\pr}{\operatorname{pr}}

\setlength{\emergencystretch}{2em}
\newenvironment{notes}
  {\begin{itemize}\setlength{\itemsep}{0.45em}\setlength{\parsep}{0pt}}
  {\end{itemize}}

\newenvironment{steps}
  {\begin{enumerate}\setlength{\itemsep}{0.45em}\setlength{\parsep}{0pt}}
  {\end{enumerate}}

\begin{document}

% The source wording, numbering, equations, and placeholders are preserved.
% Starred headings prevent LaTeX from renumbering the original headings.

\section*{1. What EPIT means for corrosion}

\begin{notes}

    \item Pitting corrosion: the protective passive film can break down locally, allowing
    corrosion to develop at these sites and form pits.

    \item Pitting potential, EPIT: the reported electrochemical potential at which
    passive-film breakdown or sustained pitting is identified during a test. It describes a
    breakdown threshold.

    \item Under comparable conditions, a more positive EPIT indicates greater resistance to
    pit initiation: a higher applied potential is needed to initiate sustained pitting.

    \item EPIT depends on alloy composition and material state, together with temperature,
    chloride concentration, pH, and the experimental procedure.

\end{notes}

\section*{2. The dataset and prediction task}

\begin{notes}

    \item The data come from Electrochemical Metrics for Corrosion Resistant Alloys,
    released on Figshare and described by Nyby et al. in Scientific Data (2021).

    \item Each row links an alloy composition and reported test conditions to a measured
    EPIT value.

    \item We predict the reported average pitting potential, “Epit, mV (SCE) Avg.”,
    expressed in millivolts relative to the saturated calomel electrode.

    \item 760 rows contain a usable numeric EPIT value and form the dataset used here.

\end{notes}

\section*{3. What information does the model receive?}

\begin{notes}

    \item Each row has 21 input columns: 17 composition variables, three environmental
    variables, and one test-method variable.

    \item The composition columns contain the weight percentages of Fe, Cr, Ni, Mo, W, Nb,
    Al, V, Ta, Re, Ce, Ti, Co, B, Mg, Y, and Gd, kept in that order.

    \item The environmental inputs are temperature in °C, chloride concentration in mol/L,
    and pH.

    \item The test-method input identifies the reported experimental method or protocol.

    \item The dataset includes Fe alloys, Ni–Cr–Mo alloys, Al alloys, high-entropy alloys,
    and other alloys.

\end{notes}

\section*{4. Keeping similar compositions together}

\begin{notes}

    \item Several rows can describe the same or very similar alloy compositions, measured
    under different conditions. We keep these rows together so that closely related
    compositions do not appear on opposite sides of an evaluation split.

    \item We compare two compositions by adding the absolute differences across their 17
    element amounts:
    \[
    d(i,j)=\sum_{e=1}^{17}|x_{i,e}-x_{j,e}|.
    \]
    For this calculation, missing amounts are treated as zero and values are rounded to 0.01
    wt.\%.

    \item Two rows are linked when their total composition difference is at most 1.0 wt.\%.
    All rows connected through these links belong to the same group.

    \item Group membership depends only on composition. Different temperatures, chloride
    concentrations, pH values, or test methods do not separate rows belonging to the same
    composition group.

    \item When assigning whole groups to partitions, we also aim to balance properties such
    as environmental conditions, alloy classes, test methods, and EPIT ranges.

\end{notes}

\section*{5. Development data, final-test data, and folds}

\begin{notes}

    \item We divide the dataset into 608 development rows (80\%) and 152 final-test rows
    (20\%), keeping every composition group together.

    \item The development data support method development and model selection. The
    final-test data assess the selected model.

    \item We divide the development data into five fixed folds, each containing 121–122
    rows. Composition groups also remain intact across these folds.

    \item For one development evaluation, one fold supplies the validation rows to predict,
    while the other four provide 486–487 labeled context rows.

    \item We repeat this five times, so every development row takes a turn as a validation
    query.

    \item For final evaluation, all 608 development rows provide labeled context, and the
    model predicts EPIT for the 152 final-test rows.

\end{notes}

\section*{2. How TabICL learns and makes predictions}

\subsection*{2.1. The prediction task and notation}

\begin{notes}

    \item Labeled context, \((x,y)\): the examples available to the model. Each \(x_i\)
    contains the inputs for one observed row, and \(y_i\) is its known response. Here, the
    inputs describe the alloy and test conditions, and the response is EPIT.

    \item Query, \(x^\ast\): a new row whose inputs are known but whose response,
    \(y^\ast\), we want to predict.

    \item Task mechanism, \(\theta\): one possible mechanism for generating a synthetic
    task, including the relationship between inputs and responses and any noise.

    \item Transformer weights, \(\phi\): the parameters TabICL learns during pretraining.
    One set of weights learns across many tasks generated from different mechanisms
    \(\theta\).

\end{notes}

\subsection*{2.2. How this differs from the usual machine-learning workflow}

\begin{notes}

    \item The usual workflow is fit, then predict. For a supervised neural network or a
    model such as CatBoost, labeled examples are used to fit the model’s parameters. The
    fitted model then predicts responses for new rows:
    \[
    (x,y)\xrightarrow{\text{fitting}}\text{fitted model}, \qquad (\text{fitted
    model},x^\ast)\longrightarrow\widehat y^\ast.
    \]

    \item A new prediction task normally requires fitting or adapting a model to its labeled
    examples.

    \item TabICL learns how to use examples during pretraining. When applied to a new task,
    it receives the labeled context and query together:
    \[
    (x,y,x^\ast) \xrightarrow{\text{pretrained TabICL}} q_\phi(y^\ast\mid x^\ast,x,y).
    \]
    Here, \(q_\phi\) represents the model’s predicted distribution of possible responses.

    \item The context guides the prediction while the transformer weights stay fixed.
    Changing the supplied examples can change the prediction without another round of
    parameter fitting. This is called in-context learning.

\end{notes}

\subsection*{2.3. TabICL and Bayesian prediction}

TabICL uses context to make predictions that approximate Bayesian prediction. It learns
this ability through synthetic training, rather than performing explicit Bayesian
calculations for each new task.

\begin{notes}

    \item A typical supervised model, such as CatBoost, is fitted to a particular labeled
    dataset. Those examples shape its fitted parameters. For a new prediction task, you
    generally fit or adapt a model again.

    \item TabICL trains one transformer across many synthetic tasks with different
    relationships. This teaches it to use the labeled context supplied for the current
    task. A new context can change its predictions while its trained weights stay fixed.

\end{notes}

To succeed, it must learn to look at the examples in the current context, work out what
relationship they suggest, and make predictions accordingly. This ability to use context
is learned during training.

The training tasks come from a synthetic-data generator. It follows a configured recipe.
In CorrPFN, this recipe includes informed and generic tasks, MLP and tree generators,
composition perturbations, and corrosion rules.

The recipe defines which \textbf{mechanisms} can be sampled and how often. This
distribution over mechanisms is the \textbf{prior}, written \(p(\theta)\).

A mechanism is one particular way of generating targets from inputs, including its noise
assumptions. For example, it could contain a particular random MLP and a selected
corrosion rule. The MLP’s particular weights matter. “Use an MLP” alone does not specify
the complete relationship.

The symbol \textbf{\(\theta\) represents the mechanism itself}. It is not a separate step
between the prior and the mechanism:

\[
\boxed{
\text{Prior}
\;\longrightarrow\;
\text{sampled mechanism }\theta
\;\longrightarrow\;
\text{synthetic dataset }D
}
\]

In sampling notation,

\[
\theta\sim p(\theta),
\qquad
D\sim p(D\mid\theta).
\]

The first expression selects the mechanism. The second generates observations under that
mechanism. The resulting dataset \(D\) contains input rows and their synthetic targets. A
training batch groups several such datasets.

One prior can therefore produce many mechanisms and many datasets. We do not define a new
prior for each table.

\textbf{Random generation does not remove prior probabilities. The sampling rules define
them.} In the selected CorrPFN configuration,

\[
P(\text{informed})=0.75,
\qquad
P(\text{generic})=0.25.
\]

These settings give informed mechanisms three times as much starting probability as
generic mechanisms, taken collectively. They express our assumption that informed
relationships deserve more emphasis for the tasks we want to solve.

This is a modeling choice. It is not a measured frequency of mechanisms occurring in
nature. Corrosion knowledge and development results help us choose the recipe.

Within the informed branch,

\[
P(\text{MLP}\mid\text{informed})=0.50.
\]

Therefore,

\[
P(\text{informed and MLP})
=
0.75\times0.50
=
0.375.
\]

This probability belongs to the whole category of informed MLP mechanisms. Further random
choices select a particular mechanism within that category. For continuously sampled
settings, \(p(\theta)\) describes a probability density.

For training, each synthetic dataset is divided into context and query rows. Write the
labeled context as

\[
D_{\mathrm{ctx}}
=
(X_{\mathrm{ctx}},y_{\mathrm{ctx}}).
\]

Here, \(X_{\mathrm{ctx}}\) contains the context inputs. The corresponding targets are
\(y_{\mathrm{ctx}}\).

For one query row, \(x^\ast\) denotes its inputs and \(y^\ast\) its target. The
transformer receives the context inputs, context targets, and query inputs. It does
\textbf{not} receive the query target.

For compact notation, call all this visible information \(O\):

\[
O=(D_{\mathrm{ctx}},x^\ast).
\]

The synthetic generator also supplies \(y^\ast\). The training procedure uses this
\textbf{hidden query target} to score the prediction. The context targets provide
examples. The query targets provide the prediction errors used for learning.

We train \textbf{one shared transformer across these tasks}. Its learned weights are
denoted by \(\phi\). Errors from different tasks update these same weights. We do not
train a separate prediction model for each synthetic dataset.

The transformer is not told which mechanism \(\theta\) produced the task. It must use the
context to make its prediction.

Bayesian reasoning describes how we could predict by explicitly considering the possible
mechanisms.

For each mechanism, the \textbf{likelihood} measures how plausibly it could produce the
visible information:

\[
L(\theta)=p(O\mid\theta).
\]

This includes the supplied inputs and observed context targets. It excludes the hidden
query target.

A mechanism receives greater support when it explains the observations well, accounting
for its noise assumptions. This is not similarity between a saved synthetic table and the
current context. It is an assessment of the observations under a possible mechanism.

The \textbf{posterior} combines this evidence with the prior:

\[
p(\theta\mid O)
=
\frac{p(O\mid\theta)\,p(\theta)}{p(O)}.
\]

The prior supplies the starting weight. The likelihood supplies evidence from the
observations. The denominator normalizes the updated weights.

This explains why the 75\% informed probability matters. Before seeing the context, the
informed branch has more support overall. If both branches explain the visible information
equally well, it retains that advantage. If the generic branch explains the observations
much better, the evidence can reverse the preference.

The prior therefore influences the prediction without fixing the final weights in advance.
\textbf{A 75\% informed-task probability does not force every prediction to contain a 75\%
informed contribution.}

Each mechanism already defines how targets arise from inputs. It can therefore predict the
query target. No additional model needs to be trained after assessing the mechanisms.

Bayesian prediction combines their predictive distributions using the posterior weights:

\[
p(y^\ast\mid O)
=
\int
p(y^\ast\mid O,\theta)\,
p(\theta\mid O)\,d\theta.
\]

The first factor is the prediction under one mechanism. The second is that mechanism’s
posterior weight. The integral combines their contributions across possible mechanisms.

Mechanisms supported more strongly by the observations contribute more. Several mechanisms
may explain the context but disagree about the query. Combining their distributions
retains this uncertainty.

\textbf{TabICL learns to approximate the resulting predictive distribution directly:}

\[
q_\phi(y^\ast\mid O)
\approx
p(y^\ast\mid O).
\]

Here, \(q_\phi\) is the distribution produced by the trained transformer. The distribution
on the right is the Bayesian prediction implied by the generator and visible information.

The implementation does not explicitly calculate the likelihoods and posterior weights
above. Instead, it trains on synthetic query targets. In CorrPFN, CRPS scores the
predicted distribution against each hidden query target. These scores drive updates to the
transformer weights \(\phi\).

The training targets are therefore the synthetic query values \(y^\ast\). They are not
calculated likelihoods or posterior weights.

The generator’s sampling probabilities influence what the transformer learns. With 75\%
informed sampling, informed tasks occur more frequently during training. This gives their
relationships greater representation in the training experience.

The context provides the other part of the learning problem. Across tasks, the transformer
encounters different visible examples and corresponding query targets. Distributional
training rewards it for predicting those query targets appropriately. Under ideal learning
conditions, it targets the same conditional distribution as the Bayesian calculation. In
practice, the result is an approximation.

When we supply real corrosion data, the transformer’s weights stay fixed. Its training has
already shaped which relationships it expects. The measured context shows what
relationship seems to apply to the current task. Both influence the predictive
distribution produced for each query.

\subsection*{2.9. Which implementation and training scope this project uses}

\begin{notes}

    \item CorrPFN and the local generic baseline use the same configured transformer
    architecture, including V2 architectural components. Both are trained from scratch.

    \item The generic MLP/tree generators and Stage 1 training framework derive from
    TabICLv1. The project adds the corrosion-informed prior and the continuous-target
    training path for EPIT.

    \item Pretrained TabICL V2 is a separate comparison model. It uses the authors’ released
    pretrained weights.

    \item Our local training uses Stage 1, with synthetic tables containing 1,024 rows.
    Later stages extend the model’s ability to use larger labeled contexts.

    \item Earlier attempts to extend training did not improve the training loss in the
    implementations tested at that time. That finding applies to those experiments and does
    not rule out improvements with newer implementations.

\end{notes}

\section*{How TabICL Learns from Synthetic Tables}

\subsubsection*{Text for presentation}

\begin{notes}

    \item An SCM transforms random starting variables into related variables.

\end{notes}

\subsubsection*{Explanation for you}

SCM means structural causal model. Mathematical functions, such as neural networks or trees,
generate new variables from earlier ones. These variables remain connected through their
dependencies.

\subsubsection*{Text for presentation}

\begin{notes}

    \item Select generated columns as features and another as the target.

\end{notes}

\subsubsection*{Explanation for you}

These columns form a synthetic dataset. Each row is one sample, with feature values and a
generated target.

\subsubsection*{Text for presentation}

\begin{notes}

    \item Provide context rows with targets and query rows without targets.

\end{notes}

\subsubsection*{Explanation for you}

The context demonstrates the relationship the model should use to predict the hidden query
targets.

\subsubsection*{Text for presentation}

\begin{notes}

    \item Represent each cell using its column, then combine features within each row.

\end{notes}

\subsubsection*{Explanation for you}

TabICL converts each cell into a vector using information from its column. Attention within
the row lets the features interact, producing a representation of the whole row.

\subsubsection*{Text for presentation}

\begin{notes}

    \item Query rows use attention to gather information from labeled context rows.

\end{notes}

\subsubsection*{Explanation for you}

Attention assigns learned weights to information from the context. This lets each query
combine useful information from the measured examples to predict its target.

\subsubsection*{Text for presentation}

\begin{notes}

    \item Compare predictions with generated targets and update the model’s parameters.

\end{notes}

\subsubsection*{Explanation for you}

The training procedure knows the hidden targets and calculates the prediction error. Updates
improve the representations, attention, and prediction output. Repeating this across tasks
with different relationships teaches TabICL how to learn from context.

We looked at the real corrosion dataset to decide what our synthetic training examples
should look like.

\begin{steps}

    \item[1.] We focused on Fe alloys and Ni–Cr–Mo alloys. For these families, we collected
    the alloy compositions and experimental conditions reported in the dataset.

    \item[2.] We used the recorded compositions as starting recipes. Each recipe is simply a
    list of element percentages. We had 298 different Fe-alloy compositions and 17 Ni–Cr–Mo
    compositions.

    \item[3.] To create a synthetic row, we randomly chose one recipe and varied its element
    amounts. For example, a chromium amount of 16 wt.\% might become 16.8 before rescaling.
    We then adjusted all amounts proportionally so the complete composition added up to 100
    wt.\%. Elements absent from the original recipe remained absent.

    \item[4.] We gave that alloy experimental conditions. We randomly selected a recorded
    combination of temperature, chloride concentration, pH, and test method. Those four
    values stayed together, but they could now be paired with a different alloy.

    \item[5.] Repeating this produced an artificial table with our 21 input columns. Each
    row described an alloy composition under particular test conditions. At this point, we
    had created the inputs; we still needed a target for each row.

    \item[6.] We generated the targets from those inputs. A random neural network or tree
    produced one response, and a corrosion-inspired formula produced another. We combined
    them into the synthetic target used for training.

\end{steps}

The central idea is: use recorded alloys and conditions to construct the inputs, then
generate artificial relationships between those inputs and the target.

The goal was to make the synthetic relationships more relevant to corrosion by including
knowledge from the literature. A random SCM can generate useful mathematical relationships,
but it has no built-in reason to associate increasing chloride with decreasing pitting
resistance: the relationship could point either way. We added corrosion rules to give part
of the generated response a physically motivated direction.

The motivation includes the role of Cr-rich passive films, contributions from Mo, and
chloride-induced film damage, supported by passive-film measurements and direct observations
of chloride attack. The precise equations and thresholds below are our modeling choices for
representing such tendencies.

For example, consider the rule that separates the environmental effects.

It uses a PREN-inspired material contribution because PREN provides a familiar way to
summarize composition-related pitting resistance. We extend this idea by explicitly
including the experimental conditions:

    \[
    y_{\mathrm{rule}} = aP-bC-cH-dI-eA.
    \]

Each letter represents a score calculated from the row’s inputs:

\begin{notes}

    \item Material protection \(P\): start from

    \[
    \mathrm{Cr}+3.25(\mathrm{Mo}+0.55\,\mathrm{W}),
    \]

using element amounts in wt.\%. More Cr, Mo, or W increases this quantity. We standardize it
and apply a smooth transformation that levels off at high values. This is PREN-inspired;
nitrogen is absent because it is outside our retained inputs.

    \item Chloride aggressiveness \(C\): take the logarithm of chloride concentration,
    standardize it, and pass it through a smooth increasing curve between zero and one. The
    logarithm represents changes across orders of magnitude. Higher chloride produces a
    larger penalty.

    \item High-temperature contribution \(H\): convert temperature into a score that rises
    smoothly from near zero to near one, with its midpoint at \(50\,^\circ\mathrm{C}\).
    Higher temperature therefore increases this penalty.

    \item Temperature–chloride interaction \(I\): multiply the temperature and chloride
    scores. This adds a penalty that becomes strongest when both are high.

    \item Acidic-pH contribution \(A\): calculate a score that increases smoothly as pH
    decreases, with its midpoint at pH 6.5. More acidic conditions therefore increase this
    penalty.

\end{notes}

We standardize each completed contribution before combining them: subtract its mean and
divide by its standard deviation. This puts them on comparable scales. During calibration,
these statistics come from the fitting rows; during synthetic generation, they come from the
generated table.

The weights \(a,b,c,d,e\) determine how strongly each contribution affects the response.
Material protection raises it; the environmental contributions lower it. The weights are
fitted from development measurements, as described below.

We proposed seven versions, each adding a particular relationship.

\begingroup
\setlength{\extrarowheight}{2pt}
\begin{longtable}{@{}>{\raggedright\arraybackslash}p{0.30\textwidth}>{\raggedright\arraybackslash}p{\dimexpr0.70\textwidth-2\tabcolsep\relax}@{}}
\hline
\textbf{Rule} & \textbf{Relationship it contributes} \\[0.45em]
\hline
Linear PREN-inspired & Cr, Mo, and W increase the material-resistance score. Environmental aggressiveness and material-dependent chloride effects lower the response. \\[0.45em]
Cr–Mo/W synergy & Adds a joint Cr–Mo/W contribution: the benefit depends on their combination. It also includes an additional penalty when temperature and chloride are both elevated. \\[0.45em]
Threshold and saturation & Makes the chromium contribution rise most strongly around 12 wt.\% and then level off. Additional Mo/W gives progressively smaller increases. This introduces a nonlinear composition response. \\[0.45em]
Separate environmental effects & Gives chloride, high temperature, temperature–chloride interaction, and acidic pH their own weights, allowing their relative importance to be fitted separately. \\[0.45em]
Coupled breakdown & Combines environmental attack and material resistance inside one smooth transition. A stronger material score shifts that transition toward more aggressive conditions. \\[0.45em]
Fe/Ni chromium threshold & Uses different chromium transition points: 12 wt.\% for Fe alloys and 15 wt.\% for Ni–Cr–Mo alloys. This allows the same chromium amount to contribute differently in the two families. \\[0.45em]
Method-aware & Extends the synergy rule with a method-dependent shift, allowing the reported response to vary with the experimental procedure. \\[0.45em]
\hline
\end{longtable}
\endgroup

For example, the temperature penalty is motivated by measurements showing decreasing pitting
potential with increasing temperature in tested austenitic stainless steels.
Temperature-dependence study

These rules express candidate tendencies. Their chosen thresholds and pH dependence are
assumptions we evaluated against our data.

We fitted their relative weights using real measurements.

We used the 452 Fe-alloy and Ni–Cr–Mo rows in the development partition. For each formula,
we adjusted its weights to make its scores match the ranking of measured EPIT values:
measurements with higher EPIT should generally receive higher rule scores.

The mathematical terms and internal thresholds stayed fixed. The fitted weights were
nonnegative and summed to one, preserving the intended directions. Fitting also discouraged
large departures from predefined starting weights.

We evaluated each rule on rows excluded from its fit.

For each of the five development folds, we fitted the rule and its preprocessing on the
other four folds, then calculated its scores on the excluded fold. We measured agreement
with the observed EPIT ordering using Spearman correlation, then averaged the five
correlations.

A score of 1 means perfect agreement in ordering; 0 means no rank correlation. The recorded
results were:

\begingroup
\setlength{\extrarowheight}{2pt}
\begin{longtable}{@{}>{\raggedright\arraybackslash}p{0.55\textwidth}>{\raggedright\arraybackslash}p{\dimexpr0.45\textwidth-2\tabcolsep\relax}@{}}
\hline
\textbf{Rule} & \textbf{Mean validation Spearman} \\[0.45em]
\hline
Linear PREN-inspired & 0.5174 \\[0.45em]
Cr–Mo/W synergy & 0.5583 \\[0.45em]
Threshold and saturation & 0.5181 \\[0.45em]
Separate environmental effects & 0.5850 \\[0.45em]
Coupled breakdown & 0.5525 \\[0.45em]
Fe/Ni chromium threshold & 0.5252 \\[0.45em]
Method-aware & 0.5999 \\[0.45em]
\hline
\end{longtable}
\endgroup

The method-aware rule gave the strongest ordering agreement. These scores assess the
formulas directly; transformer performance is evaluated later.

We used both the fitted weights and the validation scores during synthetic generation.

After validation, we refitted each rule on all 452 eligible development rows to obtain its
final weights. We retained all seven rules and sampled them in proportion to their
validation scores: approximately 15.6\% for the method-aware rule and 13.4\% for linear
PREN, for example.

For each informed task, one selected rule generates a response across the synthetic table.
We standardize that response and the SCM response separately, combine them, and standardize
again. This adds a corrosion-informed contribution while retaining variation from the random
generator; it does not enforce the expected physical direction in every combined task.

For the method-aware rule, the fitted importance of the method contribution transfers to
generation, while fresh method-specific shifts are sampled for synthetic tasks.

Once the alloy inputs and corrosion rules were defined, we still had to decide what mixture
of synthetic tasks the model should practise on. Four settings control this.

\begin{steps}

    \item[1.] The proportion of corrosion-informed tasks, \(\rho\).

This determines whether a generated task follows our corrosion construction or the generic
SCM generator.

Increasing \(\rho\) gives the model more practice with alloy compositions, test conditions,
and corrosion-related responses. Decreasing it gives more practice with general mathematical
relationships.

The selected setting was \(\rho=0.75\): an informed task is chosen with 75\% probability and
a generic task with 25\% probability.

    \item[2.] The amount of composition variation, \(\tau\).

Within an informed task, we start from recorded alloy compositions and randomly multiply
their positive element amounts by factors around one. This creates variations around the
recorded recipes.

A small \(\tau\) keeps those factors close to one. A larger \(\tau\) allows stronger
changes. Elements absent from the starting composition remain absent, and the complete
composition is rescaled to 100 wt.\% afterward.

The selected setting was \(\tau\approx0.103\). It controls the spread of the random
multipliers; it is not an amount added to every element’s weight percentage.

    \item[3.] The balance between neural-network and tree generators.

Once the physical inputs are constructed, a random MLP or tree generates the SCM part of the
response.

An MLP creates relationships through successive neural-network transformations. A tree
creates relationships through branching decisions based on input values. Changing their
sampling probabilities changes the kinds of relationships the model encounters.

The selected setting used equal probabilities for MLPs and trees within informed tasks. The
generic branch retained its existing 70\% MLP and 30\% tree mixture.

    \item[4.] The strength of the corrosion-rule contribution, \(\lambda\).

Each informed task has two responses calculated from its inputs: one from the random SCM and
one from a selected corrosion rule. After standardizing them separately, we combine them:

    \[
    y_{\mathrm{combined}} = (1-\lambda)y_{\mathrm{SCM}} + \lambda y_{\mathrm{rule}}.
    \]

At \(\lambda=0\), the target comes entirely from the SCM. At \(\lambda=1\), it comes
entirely from the corrosion rule. Intermediate values retain random relationships while
adding the chosen corrosion tendencies. The combined response is then standardized again.

The selected setting was \(\lambda\approx0.205\): approximately 80\% SCM weight and 20\%
corrosion-rule weight.

\end{steps}

The distinction between the two mixture settings is useful: \(\rho\) controls how often an
informed task is generated; \(\lambda\) controls the target inside that informed task. The
term weights fitted earlier control the contributions within the selected corrosion rule.

We used Optuna to select these four settings by comparing short transformer-training runs on
the development folds, with the calibrated rules kept fixed.

The best completed configuration at selection was Trial 56:

\begingroup
\setlength{\extrarowheight}{2pt}
\begin{longtable}{@{}>{\raggedright\arraybackslash}p{0.45\textwidth}>{\raggedright\arraybackslash}p{\dimexpr0.55\textwidth-2\tabcolsep\relax}@{}}
\hline
\textbf{Setting} & \textbf{Selected value} \\[0.45em]
\hline
Informed-task proportion \(\rho\) & 0.75 — 75\% informed, 25\% generic \\[0.45em]
Composition perturbation strength \(\tau\) & 0.1028 \\[0.45em]
MLP probability within informed tasks & 0.50 — equal MLP/tree probabilities \\[0.45em]
Corrosion-rule mixture weight \(\lambda\) & 0.2048 — 20.48\% rule weight, 79.52\% SCM weight \\[0.45em]
\hline
\end{longtable}
\endgroup

After 1,000 training steps, this configuration achieved a mean development-fold Spearman
correlation of 0.6903, the criterion used to select it.

It used the 21 original input columns, with Magpie descriptors disabled.

\end{document}

